\clearpage\fancyhead[L]{\small INSTRUMENTS.}\hypertarget{native-370-lower}{}\label{native-370-lower}
\begin{center}\textbf{INSTRUMENTS.}\par\vspace{4pt}\hrule\vspace{4pt}\fontsize{10}{12}\selectfont Göttingische gelehrte Anzeigen. No. 75. pp. 745\;..\;752. 10 May 1813.\par\vspace{4pt}\hrule\vspace{4pt}\end{center}

\par\indent At the meeting of the Royal Society of Sciences on 3 April, at which Professor \textsc{Eichhorn} gave the lecture, Professor \textsc{Gauss} reported on two new astronomical instruments. Our observatory acquired them several months ago from the workshop of the renowned instrument maker Mr. \textsc{Reichenbach}, a councillor of the saltworks in Munich. Professor \textsc{Gauss} also presented the first observations made with them.

\par\indent The \emph{first} is a repeating circle with two telescopes and a diameter of 12 inches. In this work of art one admires the fineness and accuracy of the graduations, the almost incredible sensitivity of the spirit levels and the perfection of the telescopes as much as the precision and beauty with which each individual part is finished, the ease of all its many motions and the various ingenious devices by which the gifted maker has made the observer’s work easier.

\par\indent The main circle is graduated directly on inlaid silver at intervals of 5 minutes. Each of the four verniers in turn divides one such interval into 75 parts, or steps of 4 seconds. The verniers are spaced at right angles and lie in the plane of the graduated arc on a complete second circle. This circle moves so perfectly concentrically inside the first that, although the unaided eye can see no gap at all, \SourceLeaf{371}it nowhere touches the first circle and therefore moves with the greatest ease. The gap is also invisible through the deliberately slightly tilted microscopes. The vernier marks seem to meet those of the graduated arc directly. This greatly improves the accuracy of reading, as do suitable lighting and the extraordinary fineness and uniformity of the graduation marks. One can readily estimate 2 seconds and, if necessary, even single seconds. The four verniers never differ by more than a couple of seconds. This demonstrates both the absence of any perceptible eccentricity and the unsurpassed accuracy of the graduations themselves.

\begingroup\rightskip=0pt plus 1em \par\indent The three spirit levels belonging to the circle are exceptionally sensitive. For a tilt of one second, the main level on the rear telescope gives a deflection of more than one Paris line. The extremely slight yet uniform curvature inside the glass tube implies a radius of more than 1400 feet and could be obtained only by very skilled and delicate workmanship. To preserve this curvature, the glass tubes were sealed completely at both ends with precisely ground glass stoppers covered with a bladder membrane, without melting the ends shut.\par\endgroup

\par\indent Although the two telescopes are only 16 inches long and have an aperture of 15 lines, they are of quite exceptional quality, perfectly achromatic, and their extraordinary precision allows high magnification. Sighting on the fine spider threads drawn across the focus can therefore be done with great precision. The front telescope has a prismatic eyepiece, so that when measuring altitude the eye always looks along horizontal rays. Stars high in the sky are thus as easy to observe as low ones. The loss of light is almost imperceptible.

\par\indent All the screws used for fine adjustment are made with the utmost precision. They respond properly and immediately to the lightest touch, and special brake screws prevent any future backlash. Their threads are so fine that more than 100 turns fit in an inch.

\par\indent We pass over several features peculiar to \textsc{Reichenbach}’s circles, as ingenious as they are practical, since this is not the place for a full description of the instrument. Among them are ways of setting the plane of the circle exactly vertical and of finding objects easily and bringing them into the field of view, among other things.

\par\indent Several circumstances delayed the start of astronomical observations with this circle until the middle of March. Since then Professor \textsc{Gauss} has made observations with the help of Professor \textsc{Harding}, who took care of setting the level. As a sample we give the results for the altitude of the pole obtained by observing the lower culminations of Polaris. Their close agreement is the best evidence of the excellence of the instrument.

\Needspace{13\baselineskip}\begin{center}Pole altitude of the Göttingen Observatory from observations of Polaris at lower culmination\end{center}
\begin{center}\fontsize{10}{12}\selectfont\setlength{\tabcolsep}{7pt}\renewcommand{\arraystretch}{1.15}\begin{tabular}{ll|c|r}
& 1813 & \shortstack{Number of\\observations} &\\\hline
March & 20 & 10 & $51^\circ\;31'\;54''86$\\
& 22 & 18 & $55.73$\\
& 26 & 18 & $57.41$\\
& 31 & 18 & $56.25$\\
April & 3 & 18 & $56.55$\\
& 7 & 32 & $55.20$\\
& 8 & 22 & $57.33$\\
\end{tabular}\end{center}

\par\indent\SourceLeaf{372}The mean of all 136 observations is $51^\circ\,31'\,56''20$, from which the most extreme result differs by only $1''34$. A reduction of $0''16$ to the center of the observatory still has to be subtracted. The observatory’s pole altitude is therefore $51^\circ\,31'\,56''04$, only $2''$ greater than the value determined by \textsc{Tob. Meyer}. This result still depends on the declination of Polaris, for which Mr. von \textsc{Zach}’s determination was used. It may therefore change once Polaris can be observed at upper culmination, though the change will certainly be very small. Had one used for the declination the mean of the determinations of von \textsc{Zach}, \textsc{Oriani}, \textsc{Bouvard} and \textsc{Pond} (whose extremes differ by only $\frac{3}{4}$ of a second), or \textsc{Pond}’s determination alone, which agrees exactly with that mean, the pole altitude would have turned out $0''44$ smaller still.

\par\indent The \emph{second} instrument, an equally admirable masterpiece of its kind, is a repeating theodolite 8 inches in diameter. The horizontal circle is graduated directly at intervals of 10 minutes. Each of the four verniers reads to $10''$, and smaller divisions can still be estimated. Here too the verniers lie on a complete circular disk inside the graduated circle and in the same plane. The inner or vernier circle carries two supports, from which the main telescope, 12 inches long with an aperture of 13 lines, is suspended as in a transit instrument on a horizontal axis ending in two perfectly equal cylindrical steel pivots. Its line of sight, just like that of a meridian telescope, can be brought exactly perpendicular to this axis by reversing the telescope. A level hung on the axis serves to set the plane of the circle horizontal and the axis parallel to it. On this level, 1 second gives a deflection of half a line. Before touching the horizontal circle the telescope can be inclined up to $40^\circ$ above and below the horizon. The inclination is measured on a separate vertical circle 5 inches in diameter. Its vernier reads to single minutes, and half or a third of a minute can still be estimated. The second, lower telescope is exactly like the upper one, but serves only as a check that the instrument stays fixed. It is astonishing how precisely angles can be measured. with so small an instrument. Individual measurements always give angles to within a few seconds, and by repetition one can be sure of the true values to within $1''$, or at most $2''$. Especially important for astronomical use is the ease and accuracy with which observations of the Sun or the fixed stars (first-magnitude stars are easily seen in broad daylight) allow one to determine the azimuths of terrestrial objects. Once these have been determined precisely for a given station, they in turn offer a convenient way to determine the time independently of refraction. This is a great advantage, particularly in the winter months at an observatory without a transit instrument.

\par\indent It is especially gratifying that the fine telescopes on these instruments are entirely of German origin. In quality they match the best English telescopes of the same dimensions and may even surpass them. The flint glass for the objectives is made at Benedictbeuern. --- [\textsc{Gauss} appended to this report a price list of the optical instruments made at the optical institute of \textsc{Reichenbach}, \textsc{Frauenhofer} and von \textsc{Utzschneider}. The list is omitted from this edition of his works.]
